\textbf{Program 1: Graph storage, BFS, DFS, and exact-length paths.}

\textbf{File:} \texttt{graph\_traversal.cpp}.

\textbf{Chapter connection:} Adjacency list; Breadth-first search; Depth-first search.

Store an undirected edge in both lists. BFS marks on enqueue; DFS records on entry. The exact-length simple-path search uses branch-local markers and undoes a failed branch. Stored neighbour order determines the printed traversal order.

\begin{utpcompactlabfigure}[How this program runs]
\begin{labgraphic}
\begin{tikzpicture}[font=\small,execute at begin node={\hyphenpenalty=10000\exhyphenpenalty=10000}]
\node[utpstep,text width=2.50cm,minimum height=1.15cm] (s0) at (0.0,0) {\textbf{1}\\Store both directions};
\node[utpstep,text width=2.50cm,minimum height=1.15cm] (s1) at (2.95,0) {\textbf{2}\\BFS with queue};
\node[utpstep,text width=2.50cm,minimum height=1.15cm] (s2) at (5.9,0) {\textbf{3}\\DFS with call stack};
\node[utpstep,text width=2.50cm,minimum height=1.15cm] (s3) at (8.85,0) {\textbf{4}\\Backtrack exact path};
\draw[utparrow] (s0.east)--(s1.west);
\draw[utparrow] (s1.east)--(s2.west);
\draw[utparrow] (s2.east)--(s3.west);
\end{tikzpicture}
\end{labgraphic}
\end{utpcompactlabfigure}


\begin{utpsteps}\setcounter{enumi}{0}\item \textbf{Headers and initial definitions.}\end{utpsteps}

\texttt{graph\_traversal.cpp} --- lines 16--29.

\lstinputlisting[firstline=16,lastline=29]{all_experiments/Lab09/graph_traversal.cpp}


\tcbbreak

\begin{utpsteps}\setcounter{enumi}{1}\item \textbf{Implement bfs, visited, dfs.}\end{utpsteps}

\texttt{graph\_traversal.cpp} --- lines 30--62.

% Copied from chapters/09-graphs.tex (fig:bfs).
% Keep the chapter drawing, type sizes, colours, and node styles unchanged.
\begin{adslabfigure}[BFS with queue states]
\begin{adslabgraphic}
\begin{tikzpicture}[font=\small]
% graph with BFS levels
\begin{scope}
\node[utpnodedone] (A) at (0,1.5) {A};
\node[utpnodehot] (B) at (2.1,1.5) {B};
\node[utpnodehot] (C) at (0,0) {C};
\node[utpnodehot] (D) at (2.1,0) {D};
\node[utpnode] (E) at (3.9,0.75) {E};
\foreach \u/\v in {B/E,C/D,D/E,B/D}{\draw[utpedge,draw=UTPMuted!40] (\u)--(\v);}
\foreach \u/\v in {A/B,A/C,A/D}{\draw[utpedge,draw=UTPOrange,line width=1.1pt] (\u)--(\v);}
\draw[utpedge,draw=UTPGreen,line width=1.1pt] (B)--(E);
\node[utplabel,anchor=south,color=UTPGreen] at (0,2.05) {level 0};
\node[utplabel,anchor=south,color=UTPOrange] at (2.1,2.05) {level 1};
\node[utplabel,anchor=south,color=UTPBlue] at (3.9,1.35) {level 2};
\end{scope}
% queue trace
\begin{scope}[xshift=6.9cm,yshift=1.35cm]
\foreach \act/\qs [count=\i from 0] in {%
  {visit A}/{B C D},
  {pop B, find E}/{C D E},
  {pop C}/{D E},
  {pop D}/{E},
  {pop E}/{empty}}{
  \node[utplabel,anchor=east,font=\scriptsize\ttfamily] at (0,-0.56*\i) {\act};
  \node[utpcell,minimum width=2.1cm,minimum height=0.44cm,
        font=\scriptsize\ttfamily\color{UTPNavy}] at (1.35,-0.56*\i) {\qs};
}
\node[utplabel,anchor=south,color=UTPNavy] at (1.35,0.32) {\textbf{queue after step}};
\node[utplabel,anchor=north,align=center] at (0.9,-2.75) {output: A B C D E};
\end{scope}
\end{tikzpicture}
\end{adslabgraphic}
\end{adslabfigure}


\lstinputlisting[firstline=30,lastline=62]{all_experiments/Lab09/graph_traversal.cpp}


\tcbbreak

\begin{utpsteps}\setcounter{enumi}{2}\item \textbf{Implement dfsStack, visited, pathOfLength.}\end{utpsteps}

\texttt{graph\_traversal.cpp} --- lines 63--97.

\begin{utpcompactlabfigure}[DFS: returning from C resumes the suspended call at D]
\begin{labgraphic}
\begin{tikzpicture}[font=\small]
\foreach \states/\msg [count=\i from 0] in {{A}/{enter A}, {A B}/{A calls B}, {A B D}/{B calls D}, {A B D C}/{D calls C}, {A B D}/{C returns; D continues}, {A B D E}/{D calls E}} {
\node[utpcell,minimum width=4.2cm] at (0,-.75*\i) {\texttt{\states}};
\node[utplabel,anchor=west] at (3,-.75*\i) {\msg};
}
\node[utplabel,anchor=south] at (0,.4) {active recursive calls};
\node[utplabel,anchor=north,align=center] at (2,-4.4) {Output is recorded on entry: A B D C E.\\C remains visited after its call returns.};
\end{tikzpicture}
\end{labgraphic}
\end{utpcompactlabfigure}


\lstinputlisting[firstline=63,lastline=97]{all_experiments/Lab09/graph_traversal.cpp}


\tcbbreak

\begin{utpsteps}\setcounter{enumi}{3}\item \textbf{Connect the test cases in main.}\end{utpsteps}

\texttt{graph\_traversal.cpp} --- lines 98--133.

% Worked operation using existing chapter styles.
\begin{utpcompactlabfigure}[One undirected graph, two stored directions per edge]
\begin{labgraphic}
\begin{tikzpicture}[font=\small]

\node[utpnode] (A) at (0,1) {A};\node[utpnode] (B) at (2,1) {B};
\node[utpnode] (C) at (0,-1) {C};\node[utpnode] (D) at (2,-1) {D};
\node[utpnode] (E) at (4,0) {E};
\foreach \u/\w in {A/B,A/C,A/D,B/D,B/E,C/D,D/E}{\draw[utpedge] (\u)--(\w);}
\foreach \vertex/\neighbours [count=\r from 0] in {A/{B,C,D},B/{A,D,E},C/{A,D},D/{A,B,C,E},E/{B,D}}{
\node[utptag] (x\r) at (5.3,1.5-.9*\r) {\vertex};
\node[utpcell,minimum width=3.2cm] (y\r) at (7.5,1.5-.9*\r) {\neighbours};
\draw[utparrow] (x\r)--(y\r);}
\node[utplabel,anchor=north] at (1.8,-1.7) {Seven edges; fourteen neighbour entries.};
\end{tikzpicture}
\end{labgraphic}
\end{utpcompactlabfigure}


\textbf{Trace checkpoint for the supplied main.} The given seven undirected edges create fourteen adjacency entries. BFS marks discoveries so cycles do not enqueue a vertex repeatedly. DFS follows the stored neighbour order; reversing lists changes the visiting sequence without changing the reachable set. Distinguish global visited for traversal from onPath in the simple-path search, where backtracking must release the current branch's mark.

\lstinputlisting[firstline=98,lastline=133]{all_experiments/Lab09/graph_traversal.cpp}


\tcbbreak

\begin{utpsteps}\setcounter{enumi}{4}\item \textbf{Compare traversal orders and reverse neighbour order.}\end{utpsteps}

\texttt{graph\_traversal.cpp} --- lines 134--163. Continue in the same file; do not add another header or opening brace.


\lstinputlisting[firstline=134,lastline=163]{all_experiments/Lab09/graph_traversal.cpp}


\tcbbreak

\begin{utpsteps}\setcounter{enumi}{5}\item \textbf{Test exact-length paths with fresh visit markers.}\end{utpsteps}

\texttt{graph\_traversal.cpp} --- lines 164--181. Continue in the same file; do not add another header or opening brace.

% Worked lab input; chapter fonts, colours and TikZ styles retained.
\begin{adslabfigure}[Exact-length search: undo the failed branch]
\begin{adslabgraphic}
\begin{tikzpicture}[font=\small]
\node[utpnodehot,minimum size=8mm] (r0n0) at (0.00,0.00) {A};
\node[utplabel,anchor=west,align=left,text width=5.5cm] at (3.8,0.00) {Initial call: remaining=2; mark A};
\node[utpnode,minimum size=8mm] (r1n0) at (0.00,-1.35) {A};
\node[utpnodehot,minimum size=8mm] (r1n1) at (1.35,-1.35) {B};
\draw[utparrow] (r1n0)--(r1n1);
\node[utplabel,anchor=west,align=left,text width=5.5cm] at (3.8,-1.35) {Try B: push B, mark B; remaining=1};
\node[utpnode,minimum size=8mm] (r2n0) at (0.00,-2.70) {A};
\node[utpnode,minimum size=8mm] (r2n1) at (1.35,-2.70) {B};
\draw[utparrow] (r2n0)--(r2n1);
\node[utpnodehot,minimum size=8mm] (r2n2) at (2.70,-2.70) {D};
\draw[utparrow] (r2n1)--(r2n2);
\node[utplabel,anchor=west,align=left,text width=5.5cm] at (3.8,-2.70) {Try D: remaining=0, but D is not E\\Undo: pop D and clear its on-path marker};
\node[utpnode,minimum size=8mm] (r3n0) at (0.00,-4.05) {A};
\node[utpnode,minimum size=8mm] (r3n1) at (1.35,-4.05) {B};
\draw[utparrow] (r3n0)--(r3n1);
\node[utpnodehot,minimum size=8mm] (r3n2) at (2.70,-4.05) {E};
\draw[utparrow] (r3n1)--(r3n2);
\node[utplabel,anchor=west,align=left,text width=5.5cm] at (3.8,-4.05) {Try E: remaining=0 and target matches\\Return true; retain the successful path};
\end{tikzpicture}
\end{adslabgraphic}
\end{adslabfigure}


\lstinputlisting[firstline=164,lastline=181]{all_experiments/Lab09/graph_traversal.cpp}


\tcbbreak

\textbf{Build, run, and inspect graph\_traversal.}

From this program\textquotesingle s folder, run these commands in PowerShell:

\begin{lstlisting}[style=utpterminal,language={}]

g++ -std=c++17 -Wall -Wextra graph_traversal.cpp -o graph_traversal.exe

.\graph_traversal.exe

\end{lstlisting}

\textbf{Expected output excerpt.} Compare values, order, and pass/fail messages; spacing need not match.

\begin{lstlisting}[style=utpterminal,language={}]

== Adjacency list (5 vertices, 7 edges) ==
  A: B C D 
  B: A D E 
  C: A D 
  D: A B C E 
  E: B D 
  [ok]   every edge is stored in both endpoints' lists exactly once (undirected)
  [ok]   the degree sum is twice the edge count

== Traversals from A ==
  BFS (queue):          A B C D E 
  DFS (recursion):      A B D C E 
  DFS (explicit stack): A B D C E 
  [ok]   BFS gives A B C D E: all of A's neighbours before E
  [ok]   DFS gives A B D C E: it follows one branch as deep as it can

\end{lstlisting}

The complete captured run is in \texttt{labs/outputs/graph\_traversal.txt}.

\begin{tryit}

Reverse neighbour order and repeat both traversals. For the exact-length path A to E, lengths 1 and 5 must fail; 2 and 4 have simple paths. Reset markers before each new query.

\end{tryit}
